\subsection{正弦波加窄带高斯过程}
记相位 $\theta\sim U(0,2\pi)$ 的正弦波叠加窄带高斯噪声的信号表达式为
\begin{equation} \label{eq:3.6 1}
    r(t)=A\cos(\omega_ct+\theta)+n(t).
\end{equation}
其中 $n\sim N(0,\sigma_n)$, 且
\begin{equation}
    n(t)=n_c(t)\cos\omega_ct-n_s(t)\sin\omega_ct.
\end{equation}

展开式 (\ref{eq:3.6 1}), 得到
\begin{equation}
    \begin{aligned}
        r(t) & =[A\cos\theta+n_c(t)]\cos\omega_ct-[A\sin\theta+n_s(t)]\sin\omega_ct \\
             & =z_c(t)\cos\omega_ct-z_s(t)\sin\omega_ct                             \\
             & =z(t)\cos(\omega_ct+\varphi(t)).
    \end{aligned}
\end{equation}

容易得到
\begin{equation*}
    E[z_c]=A\cos\theta, E[z_s]=A\sin\theta, \sigma_c^2=\sigma_s^2=\sigma_n^2.
\end{equation*}

$f(z)$ 的概率密度函数为
\begin{equation}
    f(z)=\frac{z}{\sigma^2}\exp\left[-\frac{1}{2\sigma^2}(z^2+A^2)\right]I_0\left(\frac{Az}{\sigma^2}\right). \qquad (z\geq 0)
\end{equation}
其中 $I_0(\cdot)$ 为第一类零阶修正贝塞尔函数. 此时认为 $z$ 满足\textbf{莱斯分布}.

两种极端情况:
\begin{itemize}
    \item 信号很小, $A\rightarrow0$, $z$ 退化为瑞利分布;
    \item 信噪比很大, $\displaystyle\gamma=\frac{A^2}{2\sigma_n^2}$, $z$ 退化为高斯分布.
\end{itemize}
